\clearpage\fancyhead[L]{\small AUSZUG AUS EINEM SCHREIBEN.}\hypertarget{native-323-lower}{}\label{native-323-lower}

\begin{center}\textbf{AUSZUG AUS EINEM SCHREIBEN DES HRN. PROF. GAUSS.}\par\vspace{4pt}\hrule\vspace{4pt}\fontsize{10}{12}\selectfont Monatliche Correspondenz. Band XVIII. S. 289. 290. 1810 September.\par\vspace{4pt}\hrule\vspace{4pt}\end{center}
\begin{flushright}Göttingen, den 17. August 1810.\end{flushright}

\par\indent ---\enspace---\enspace--- Dass unser König jetzt die Fortsetzung des Baues unserer neuen Sternwarte decretirt und mit wahrhaft königl. Freigebigkeit 200000 Franken dazu bewilligt hat, werden Sie bereits wissen
\SourceLeaf{324}Hoffentlich werden Sie den Bau, wenn Sie diesen Herbst noch zu uns kommen, bereits wieder angefangen finden.

Am 8. August war ich in Münden, zum Theil in der Absicht, um den Einfluss des Fahrens auf den Gang unsers Chronometers zu finden. Ich wurde indessen vom Wetter sehr schlecht begünstigt und konnte Vormittags zwischen 10 und 11 Uhr nur einige schlechte Sonnenhöhen erhalten. Mittags war die Sonne ganz unsichtbar, und erst nach 1 Uhr erhielt ich einige zuverlässige Höhen, aus denen ich verbunden mit einigen andern gleichfalls guten, gegen 6 Uhr genommenen, die Polhöhe $51^\circ\ 25'\ 22''$ und den Längenunterschied von \emph{Göttingen} $1^{\text{m}}9^{\text{s}}3$ westlich in Zeit gefunden habe. Der Chronometer hat sich vortrefflich gehalten, und sein täglicher Gang vom 7\,..\,9. August auf $0^{\text{s}}1$, derselbe wie vorher. Da meine Sonnenhöhen so weit vom Mittage ab lagen, so will ich über die Genauigkeit meiner Bestimmung nicht eher urtheilen, als bis ich einmal Gelegenheit habe, sie zu wiederholen; doch muss ich bemerken, dass meine neun Sonnenhöhen gut unter einander harmoniren. Mein Beobachtungsplatz war der sogenannte \emph{Freitagswerder}, beim Zusammenfluss der \emph{Fulda} und \emph{Werra} gleich nordlich  vor der Stadt. Professor \textsc{Seyffer} fand 1794 auf einem vermuthlich etwas \emph{südlicher} liegenden Punkte die Polhöhe $51^\circ\ 26'\ 52''$, jedoch nur aus einer einzelnen Höhe.
\clearpage\fancyhead[L]{\small GAUSS AN BODE.}

\begin{center}\textbf{GAUSS AN BODE.}\par\vspace{4pt}\hrule\vspace{4pt}\fontsize{10}{12}\selectfont Astronomisches Jahrbuch für 1813. Seite 236\,..\,238. Berlin 1810.\par\vspace{4pt}\hrule\vspace{4pt}\end{center}
\begin{flushright}Göttingen, den 2. September 1810.\end{flushright}

\par\indent ---\enspace---\enspace---

\par\indent Die Elemente, nach welchen diese Ephemeride [der Pallas für 1810 Sept. 16 bis 1811 Juli 1 von \textsc{Gauss}] berechnet ist, gründen sich nur auf die 4 letzten Oppositionen, und werden hinreichend sein, den Planeten diesmal aufzufinden. Ich habe aber auch vor einiger Zeit diejenigen Elemente berechnet, welche sich so genau als möglich an alle 6 bisher beobachteten Oppositionen anschliessen, und welche nun am schicklichsten zur Berechnung der Störungen werden zum Grunde gelegt werden können. Letztere fangen nach gerade an, sehr sichtbar zu werden. Ich werde über diese Rechnungen nächstens ein Memoir unsrer Societät der Wissenschaften vorlegen. Dass unser Gouvernement sehr viel für die Astronomie thun will, werden Sie bereits wissen 200000 Franken sollen zu dem Bau der neuen Sternwarte verwandt werden. Hr. Prof. \textsc{Harding} hat eine Gratification von 4000 Franken zu einer Reise nach Paris erhalten, um dort die Lücken der Histoire céleste auszufüllen und seine Sternkarten zu completiren.

Die mir neulich erzeigte Ehre der Aufnahme in Ihre Akademie der Wissenschaften ist mir sehr schätzbar.
\clearpage\fancyhead[L]{\small NEUE PALLASELEMENTE.}\SourceLeaf{325}

\begin{center}\textbf{NEUE PALLASELEMENTE.}\par\vspace{4pt}\hrule\vspace{4pt}\fontsize{10}{12}\selectfont Monatliche Correspondenz. Band XXII. S. 400\,..\,403. 1810 October.\par\vspace{4pt}\hrule\vspace{4pt}\end{center}

\par\indent Da \emph{Pallas} von den neuen Planeten zuerst sichtbar und wegen ihrer grossen Erdnähe eine bedeutende Lichtstärke haben wird, die für unsere nördlichen Zonen nur vielleicht durch ihre starke südliche Abweichung gemindert werden dürfte, so eilen wir, unsern Lesern die kürzlich von Herrn Prof. \textsc{Gauss} berechnete Ephemeride ihres Laufes hier mitzutheilen. Mit Anfang November wird sich der Planet wahrscheinlich schon im Meridian beobachten lassen.

Die Elemente, auf denen die nachfolgende Ephemeride beruht, wurden aus den Oppositionen von 1805, 7, 8 und 9 hergeleitet, und sind folgende:

\Needspace{14\baselineskip}\begin{center}\emph{Göttinger Meridian.}\par\medskip\fontsize{10}{12}\selectfont\setlength{\tabcolsep}{4pt}\renewcommand{\arraystretch}{1.15}\begin{tabular}{llr}
Mittlere Länge 1803 & .\;.\;. & $221^\circ\;23'\;24''6$\\
mittlere tägliche Bewegung &  & $770''9365$\\
1803 Perihelium & .\;.\;.\;. & $120^\circ\;58'\;4''8$\\
--- $\AscendingNode$ & .\;.\;.\;.\;.\;. & $172\;27\;52.4$\\
Neigung & .\;.\;,\;.\;.\;.\;. & $34^\circ\;36'\;49''4$\\
Excentricität & .\;.\;.\;.\;. & $0.2446335$\\
Log. halbe Axe & .\;.\;.\;. & $0.4420473$\\
\end{tabular}\end{center}

\par\indent [Der Lauf der Pallas 1810 und 1811 ist für Mitternacht in Göttingen von vier zu vier Tagen angegeben, AR. und Decl. in Minuten log Dist. von der Erde in vier Decimalen]
